% viewpoint_temporal_coexistence.tex
% Why You Must Wait: Temporal Coexistence, the Quantum Eraser,
% and the Arrow of Time as Recovery Failure
% Author: Sheng-Kai Huang
% Target: Foundations of Physics (viewpoint/perspective)
% Compile: pdflatex viewpoint_temporal_coexistence.tex

\documentclass[12pt]{article}

\usepackage{amsmath}
\usepackage{amssymb}
\usepackage{amsfonts}
\usepackage{amsthm}
\usepackage[margin=1in]{geometry}
\usepackage{hyperref}
\usepackage{bm}

\hypersetup{
  colorlinks=true,
  linkcolor=blue,
  citecolor=blue,
  urlcolor=blue
}

\newtheorem{theorem}{Theorem}
\newtheorem{corollary}[theorem]{Corollary}
\newtheorem{proposition}[theorem]{Proposition}
\theoremstyle{definition}
\newtheorem{definition}[theorem]{Definition}
\theoremstyle{remark}
\newtheorem*{remark}{Remark}

% Convenience macros (consistent with Paper 1)
\newcommand{\kB}{k_{\mathrm{B}}}
\newcommand{\Tr}{\mathrm{Tr}}
\newcommand{\id}{\mathrm{id}}
\newcommand{\Petz}{\mathcal{R}}
\newcommand{\chan}{\mathcal{N}}
\newcommand{\hilb}{\mathcal{H}}
\newcommand{\ket}[1]{|#1\rangle}
\newcommand{\bra}[1]{\langle#1|}
\newcommand{\braket}[2]{\langle#1|#2\rangle}
\newcommand{\op}[2]{|#1\rangle\langle#2|}
\newcommand{\abs}[1]{\left|#1\right|}
\newcommand{\tauasym}{\tau}

\begin{document}

\title{\textbf{Why You Must Wait: Temporal Coexistence, the Quantum Eraser,\\
and the Arrow of Time as Recovery Failure}}

\author{Sheng-Kai Huang\\
\textit{Independent Researcher}\\
\texttt{akai@fawstudio.com}}

\date{March 2026}

\maketitle

%=======================================================================
\begin{abstract}
%=======================================================================

The delayed-choice quantum eraser demonstrates that
interference can be recovered from an apparently decohered signal---but
only through coincidence counting with a spacelike-separated idler.
The standard explanation correctly invokes entanglement and
complementarity, yet leaves a question unanswered: \emph{what physical
feature of the experimenter's situation forces the ``waiting''?}
We argue that this waiting is the operational signature of a nonzero
temporal asymmetry parameter $\tauasym = 1 - F$, where $F$ is the
fidelity of Petz recovery---the unique retrodiction map satisfying
Bayesian consistency.  In a closed bipartite system (the entangled
signal-idler pair), $\tauasym = 0$: forward evolution and backward
retrodiction are informationally equivalent, and there is no operational
distinction between past and future states.  For the experimenter,
who is external to the entangled pair and must access the idler through
a classical communication channel, $\tauasym > 0$, bounded below by
the entropy production $\Sigma$ via $\tauasym \leq 1 - \exp(-\Sigma/2)$.
The no-signaling theorem provides the constraint forbidding
superluminal information transfer; $\tauasym > 0$ provides the
mechanism that creates the subjective experience of temporal delay.
We connect this framework to Aharonov's two-state vector formalism,
Price's retrocausality program, Zurek's decoherence, the
Page-Wootters mechanism, and the Connes-Rovelli thermal time hypothesis.
The resulting picture suggests that the arrow of time is not a
fundamental feature of dynamics but the operational cost of being a
subsystem in an open environment.

\end{abstract}

%=======================================================================
\section{The Quantum Eraser Paradox}
\label{sec:eraser}
%=======================================================================

Few experiments in the foundations of quantum mechanics have generated
as much confusion---both genuine and pedagogical---as the quantum eraser.
The original proposal by Scully and Dr\"uhl~\cite{Scully1982}
and its delayed-choice realization by Kim \emph{et al.}~\cite{Kim2000}
appear to demonstrate something deeply unsettling: a measurement
performed on one particle can ``restore'' interference in a distant
partner, even when the which-path information has already been
recorded.  In the delayed-choice variant, the restoration appears to
reach backward in time, since the idler measurement occurs after the
signal has already been detected.

Let us recall the essential structure.  A spontaneous parametric
down-conversion source produces entangled signal-idler pairs in the state
\begin{equation}
\ket{\Psi}_{SI} = \frac{1}{\sqrt{2}}
\big(\ket{A}\ket{d_A} + \ket{B}\ket{d_B}\big),
\label{eq:entangled}
\end{equation}
where $\ket{A}$ and $\ket{B}$ are the two path states of the
signal photon, and $\ket{d_A}$, $\ket{d_B}$ are orthogonal
states of the idler that serve as which-path markers.
When we trace out the idler---corresponding to looking at the signal
alone---the signal's reduced state is
\begin{equation}
\rho_S = \Tr_I\big(\op{\Psi}{\Psi}_{SI}\big)
= \frac{1}{2}\big(\op{A}{A} + \op{B}{B}\big),
\label{eq:decohered}
\end{equation}
which is a fully decohered mixture.  No interference is visible.
This much is straightforward quantum mechanics.

The ``eraser'' step proceeds as follows.  Instead of measuring
the idler in the which-path basis $\{\ket{d_A}, \ket{d_B}\}$,
we measure it in the complementary basis
$\ket{\pm} = (\ket{d_A} \pm \ket{d_B})/\sqrt{2}$.
Conditioned on obtaining outcome $\ket{+}$, the signal is projected
into $(\ket{A} + \ket{B})/\sqrt{2}$---a coherent superposition
that exhibits interference.  Conditioned on $\ket{-}$, the signal
is in $(\ket{A} - \ket{B})/\sqrt{2}$, which also exhibits
interference but with a $\pi$ phase shift.  The total signal
statistics (summed over both outcomes) remain the incoherent
mixture~(\ref{eq:decohered}); interference appears only in the
subensembles selected by the idler outcome.

The standard textbook explanation is essentially correct: the
correlations are present in the entangled state from the beginning,
and coincidence counting merely reveals pre-existing correlations.
No information travels faster than light.  Nothing changes in the
signal's reduced state.  The no-signaling theorem~\cite{Ghirardi1980}
guarantees that the marginal statistics at the signal detector are
independent of what measurement is performed on the idler.

And yet, something remains unexplained.  The standard account tells us
\emph{that} we must perform coincidence counting to see interference,
but it does not tell us \emph{why}.  It tells us the constraint
(no-signaling), but not the mechanism.  Consider the following
sharpened version of the question:

\begin{quote}
\emph{The entangled pair, considered as a closed bipartite system,
contains all the information needed to reconstruct interference.
This information exists ``all at once'' in the joint state
$\ket{\Psi}_{SI}$.  Why, then, does the experimenter---who is
also a physical system---experience a temporal delay in accessing it?
What physical property of the experimenter's relation to the
entangled pair creates the ``waiting''?}
\end{quote}

This is not a question about no-signaling.  No-signaling tells us that
the signal's \emph{marginal} statistics cannot depend on the idler
measurement---a kinematic constraint rooted in the tensor product
structure of Hilbert space.  But the question of why the experimenter
must perform coincidence counting---why the relevant correlations are
not directly accessible from the signal alone---is a question about
the experimenter's \emph{informational relationship} to the entangled pair.

Our answer, which we develop in the following sections, is this:
the experimenter must wait because the experimenter is \emph{external}
to the closed system that contains the coherence.  Externality means
that the experimenter's access to the idler is mediated by an
irreversible channel (classical communication, detector readout,
data storage), which introduces entropy production $\Sigma > 0$ and
therefore a nonzero temporal asymmetry $\tauasym > 0$.  The
``waiting'' is the operational signature of this nonzero $\tauasym$.

In a precise sense that we will make explicit, the entangled pair
itself does not experience waiting.  For the signal-idler system
considered as a closed bipartite system, $\tauasym = 0$: forward
evolution and backward retrodiction are informationally equivalent.
The past state (before which-path marking) and the future state
(after marking) are connected by a perfectly reversible map.
It is only when we introduce a subsystem boundary---when we ask
what the signal reveals about the idler \emph{from the perspective
of an external observer}---that temporal asymmetry appears.

The structure of this paper is as follows.
Section~\ref{sec:tau} introduces the temporal asymmetry parameter
$\tauasym$ through the Petz recovery map, citing known results without
re-derivation.
Section~\ref{sec:coexistence} develops the central interpretive
claim: that $\tauasym = 0$ implies the operational coexistence of
temporal perspectives, and connects this to existing frameworks in
the foundations of physics.
Section~\ref{sec:observer} analyzes the role of the observer and
argues that the arrow of time is the cost of subsystem perspective.
Section~\ref{sec:predictions} discusses testable consequences.
Section~\ref{sec:limitations} addresses limitations, counterarguments,
and open questions.

A note on intellectual debts.  The mathematical tools we use---the Petz
recovery map, the universal recovery bound, the Crooks fluctuation
theorem---are not new.  The connection between the quantum eraser and
Petz recovery was established in Huang~\cite{Huang2026}; the Petz
map's uniqueness as a retrodiction functor is due to Parzygnat and
Buscemi~\cite{Parzygnat2023}; the universal recovery bound to Junge
\emph{et al.}~\cite{Junge2018} and Fawzi-Renner~\cite{Fawzi2015}.
What is new here is the interpretive synthesis: using $\tauasym$ to
explain \emph{why} the experimenter must wait, connecting this to the
arrow of time, and situating the result within the broader landscape
of temporal interpretation in quantum mechanics.  We follow the
convention of clearly marking known results, new synthesis, and
interpretive claims.

%=======================================================================
\section{$\tauasym$: A Measure of Temporal Asymmetry}
\label{sec:tau}
%=======================================================================

\subsection{The Petz recovery map as retrodiction}

Given a quantum channel $\chan: \mathcal{B}(\hilb_A) \to
\mathcal{B}(\hilb_B)$ and a faithful reference state $\sigma$,
the Petz recovery map~\cite{Petz1986,Petz1988} is
\begin{equation}
\Petz_{\sigma,\chan}(\rho)
= \sigma^{1/2}\,\chan^\dagger\!\big(
\chan(\sigma)^{-1/2}\,\rho\,\chan(\sigma)^{-1/2}\big)
\,\sigma^{1/2},
\label{eq:petz}
\end{equation}
where $\chan^\dagger$ is the Hilbert-Schmidt adjoint of $\chan$.
This map is completely positive and trace-preserving~\cite{Petz1986}.
Petz's sufficiency theorem~\cite{Petz1988} states that
$\Petz_{\sigma,\chan} \circ \chan = \id$ on the support of
$\sigma$ if and only if the data-processing inequality for quantum
relative entropy is saturated.

The crucial observation, due to Parzygnat and
Buscemi~\cite{Parzygnat2023}, is that the Petz map is not merely
a useful recovery tool---it is the \emph{unique} retrodiction
functor satisfying four natural axioms: (i)~it produces a valid
quantum channel; (ii)~it composes correctly under sequential channels
(Bayesian consistency, or functoriality); (iii)~it reduces to the
identity for identity channels; and (iv)~it reduces to classical
Bayes' rule on commutative subalgebras.  There is no other
construction satisfying all four.  The Petz map is to quantum
retrodiction what Bayes' rule is to classical inference---not a
choice among options, but the unique consistent prescription.

This uniqueness has a physical consequence that has not been
sufficiently appreciated.  If we accept that retrodiction---inferring
the past from the present---should satisfy Bayesian consistency,
then the Petz map is the \emph{only} way to do it.  Any departure
from the Petz map is a departure from Bayesian consistency.
The quality of Petz recovery therefore measures, without ambiguity,
how well the past can be inferred from the present.

\subsection{Definition and bounds}

We define the temporal asymmetry parameter [known from
Huang~\cite{Huang2026}]:
\begin{equation}
\tauasym = 1 - F\!\big(\rho,\,
\tilde{\Petz}_{\sigma,\chan}(\chan(\rho))\big),
\label{eq:tau_def}
\end{equation}
where $F$ is the Uhlmann fidelity and $\tilde{\Petz}$ is the
rotated Petz map of Junge \emph{et al.}~\cite{Junge2018},
which averages the standard Petz map over modular
automorphisms to handle quantum noncommutativity.
(When $[\rho, \sigma] = 0$, the rotation is trivial and
$\tilde{\Petz} = \Petz$.)

The parameter $\tauasym$ takes values in $[0,1]$ and admits a
tight chain of equivalences and inequalities.  From the universal
recovery theorem~\cite{Junge2018}:
\begin{equation}
F^2 \geq \exp(-\Delta D),
\label{eq:jrsww}
\end{equation}
where $\Delta D = D(\rho\|\sigma) - D(\chan(\rho)\|\chan(\sigma))$
is the relative entropy decrease under $\chan$.
When $\sigma$ is a Gibbs state $e^{-\beta H}/Z$, the
Kolchinsky-Wolpert decomposition~\cite{Kolchinsky2017} and the
Reeb-Wolf bound~\cite{ReebWolf2014} relate $\Delta D$ to the
entropy production $\Sigma$, yielding:
\begin{equation}
\tauasym \leq 1 - \exp(-\Sigma/2).
\label{eq:tau_bound}
\end{equation}

The central result is the equivalence chain
[assembling~\cite{Petz1988,Fawzi2015,Junge2018,Landi2021}]:
\begin{equation}
\boxed{
\tauasym = 0 \;\Longleftrightarrow\; \Sigma = 0
\;\Longleftrightarrow\; I(A;E|B) = 0
\;\Longleftrightarrow\; \text{QMC}
}
\label{eq:equiv_chain}
\end{equation}
where QMC denotes the quantum Markov chain
condition~\cite{Fawzi2015}, and $I(A;E|B)$ is the conditional
mutual information measuring information leakage to the environment
in a Stinespring dilation.  Each individual equivalence is
established in the literature; the chain states them as a unified
result.

Let us parse this chain carefully.  The condition $\tauasym = 0$
says: \emph{the past is perfectly recoverable from the present}.
The condition $\Sigma = 0$ says: \emph{no entropy is produced}.
The condition $I(A;E|B) = 0$ says: \emph{no information leaks to the
environment}.  The condition QMC says: \emph{the system-environment
state has Markov structure, so the environment carries no non-redundant
information about the system}.  These are four ways of stating the same
physical fact.

\subsection{Application to the quantum eraser}

The quantum eraser fits this framework precisely.  The entangled
state~(\ref{eq:entangled}) is a pure bipartite state
$\ket{\Psi}_{SI}$.  The which-path channel is the partial trace
$\chan = \Tr_I$, which maps the joint state to the decohered
signal~(\ref{eq:decohered}).  Since $\ket{\Psi}_{SI}$ is pure,
the partial trace saturates the data-processing inequality:
$D(\op{\Psi}{\Psi}\|\sigma) = D(\Tr_I(\op{\Psi}{\Psi})\|
\Tr_I(\sigma))$ for any reference $\sigma$ on the joint system
whose marginal on $S$ has full support.  By Petz's sufficiency
theorem, the Petz recovery is exact:
\begin{equation}
\Petz_{\sigma,\Tr_I}\big(\rho_S\big) = \op{\Psi}{\Psi}_{SI}.
\label{eq:eraser_recovery}
\end{equation}
The coherence is perfectly recoverable.  The entropy production
is zero.  No information has been lost---it has merely been
redistributed between signal and idler.  For the bipartite system,
$\tauasym = 0$.

This is a known result [Huang~\cite{Huang2026}], but its
interpretive implications have not been fully drawn out.
The statement $\tauasym = 0$ for the signal-idler pair means:
\emph{there is no temporal asymmetry internal to the pair}.
The which-path marking and the erasure are connected by a
perfectly reversible map.  The ``past'' (before marking) and
the ``future'' (after marking) are informationally equivalent.

But the experimenter is not the signal-idler pair.  The
experimenter is external.  To access the idler's measurement
outcome, the experimenter must use a classical communication
channel: a detector fires, a signal is amplified, a bit is
stored in a computer memory, and eventually the experimenter
correlates the idler outcome with the signal outcome through
coincidence counting.  Each of these steps---detection,
amplification, storage, communication---is an irreversible process
with $\Sigma > 0$.  By the bound~(\ref{eq:tau_bound}), the
experimenter's temporal asymmetry is $\tauasym_{\text{exp}} > 0$.

This is the answer to the question posed in
Section~\ref{sec:eraser}.  The experimenter must wait because
$\tauasym_{\text{exp}} > 0$.  The waiting is not a feature of
the entangled pair; it is a feature of the experimenter's
informational relationship to the pair.  The no-signaling theorem
provides the constraint that prevents the experimenter from
accessing the correlations superluminally.  The nonzero
$\tauasym_{\text{exp}}$ provides the mechanism: the classical
channel that carries the idler outcome to the experimenter
introduces entropy production, which bounds the recovery
fidelity away from unity.

\subsection{Two distinct roles: constraint versus mechanism}

It is essential to distinguish two logically independent features
of the situation, and we wish to be explicit about this distinction
because conflating them leads to confusion.

\emph{No-signaling} is a kinematic constraint.  It states that the
marginal statistics of the signal photon are independent of
whatever measurement is performed on the idler.  This follows from
the tensor product structure of the Hilbert space and the linearity
of the partial trace.  It holds regardless of whether the systems
are spacelike separated, regardless of the experimenter's knowledge,
and regardless of any dynamical details.  No-signaling tells us that
the experimenter \emph{cannot} see interference in the signal's
marginal, period.

\emph{Temporal asymmetry} ($\tauasym > 0$) is a dynamical fact
about the experimenter's channel to the entangled pair.  It tells
us that the experimenter's retrodiction about the pair's joint
state is imperfect---that the past is not perfectly recoverable
from the present, from the experimenter's perspective.  This is
what creates the subjective experience of ``waiting'': the
experimenter cannot reconstruct the joint state from the signal
alone, and must receive additional information (the idler outcome)
through a channel that takes time to traverse.

The no-signaling theorem would hold even if the experimenter could
somehow achieve $\tauasym = 0$.  But in that (physically impossible)
scenario, the experimenter would not need to ``wait'' for the
idler outcome---the joint state would be perfectly recoverable
from the signal alone, and coincidence counting would be
unnecessary.  It is precisely because the experimenter
\emph{cannot} achieve $\tauasym = 0$---because accessing the idler
requires an irreversible classical channel---that the waiting is
experienced.

%=======================================================================
\section{What ``Present and Future Coexist'' Means}
\label{sec:coexistence}
%=======================================================================

\subsection{An operational definition}

The phrase ``present and future coexist'' is loaded with
metaphysical baggage.  We wish to strip it down to its precise
operational content, state clearly what it means and what it does
not mean, and then show how it connects to existing interpretive
frameworks.

\begin{definition}[Temporal coexistence---operational]
\label{def:coexistence}
A quantum channel $\chan$ with reference state $\sigma$ satisfies
\emph{temporal coexistence} for a state $\rho$ if and only if
$\tauasym(\rho, \sigma, \chan) = 0$, i.e., the Petz recovery
achieves unit fidelity:
$F(\rho, \Petz_{\sigma,\chan}(\chan(\rho))) = 1$.
\end{definition}

In words: temporal coexistence means that the forward channel
$\chan$ and the backward retrodiction $\Petz$ are informationally
equivalent for the state $\rho$.  Any statistical prediction about
the pre-channel state $\rho$ can be equally well obtained from the
post-channel state $\chan(\rho)$, and vice versa.  The state's
informational content is time-independent: it does not matter
whether we are looking forward or backward along the channel.

Several clarifications are in order.

\emph{This is not the claim that ``the future exists.''}
Definition~\ref{def:coexistence} is an operational statement about
informational equivalence.  It is compatible with, but does not
require, ontological claims about the reality of future states.
A Bohrian instrumentalist, a many-worlds proponent, and a
retrocausalist can all accept Definition~\ref{def:coexistence}
while disagreeing about its ontological implications.  We
deliberately remain neutral on ontology.

\emph{This is not just ``unitary quantum mechanics is
reversible.''}
The claim that closed quantum systems evolve unitarily, and
that unitary evolution is reversible, is of course well known.
But Definition~\ref{def:coexistence} does more than restate
this truism in different notation.  First, it connects to the
quantum eraser \emph{specifically}, providing a quantitative
explanation for why coincidence counting is needed.
Second, $\tauasym$ is a \emph{continuous} parameter, not a
binary distinction: systems can be ``partially temporally
asymmetric'' in a precisely quantified way.
Third, the condition $\tauasym = 0$ is characterized by the
equivalence chain~(\ref{eq:equiv_chain}), linking it to
entropy production, environmental information leakage, and the
quantum Markov condition---connections that the bare statement
``unitary evolution is reversible'' does not provide.
Fourth, the bound~(\ref{eq:tau_bound}) connects temporal
asymmetry to thermodynamics through the Crooks fluctuation
theorem~\cite{Crooks1999}, giving $\tauasym$ a
thermodynamic interpretation.

\emph{Coexistence is not retrocausality.}
The condition $\tauasym = 0$ does not mean that the future
causes the past, or that information travels backward in time.
It means that there is no operational distinction between
forward and backward inference---that the channel and its Petz
reversal are equally good.  This is weaker than retrocausality
and stronger than mere reversibility.  It is, in a sense,
\emph{a-causality}: when $\tauasym = 0$, the concepts of
``cause'' and ``effect'' lose their operational grounding, because
neither temporal direction is informationally privileged.

\subsection{Connections to existing frameworks}

The $\tauasym$ framework does not emerge in isolation.  It
connects to, and in some cases provides a quantitative refinement
of, several existing interpretive programs.  We survey these
connections, noting where $\tauasym$ adds something new and where
it merely restates known positions.

\subsubsection{Aharonov's two-state vector formalism}

The two-state vector formalism (TSVF)~\cite{ABL1964,Aharonov1990}
describes a quantum system by both a forward-evolving state
$\ket{\psi(t)}$ and a backward-evolving state $\bra{\phi(t)}$,
determined respectively by pre- and post-selection.
The ABL rule~\cite{ABL1964} gives the probability of an
intermediate measurement outcome conditioned on both boundary
conditions, and this probability is time-symmetric.

In the $\tauasym$ language: \emph{the TSVF works perfectly in
the regime $\tauasym = 0$}.  When the system is closed and the
evolution is unitary, the forward and backward states carry
the same information, and the ABL rule yields symmetric
probabilities.  The TSVF's power---and its strangeness---come
from the fact that it treats both temporal directions on equal
footing, which is precisely the content of $\tauasym = 0$.

What $\tauasym$ adds is a quantitative account of when and how
the TSVF \emph{fails}.  In an open system with $\tauasym > 0$,
the backward-evolving state carries less information than the
forward-evolving state (or more precisely, the retrodiction from
the post-channel state is imperfect).  The ABL probabilities are
no longer exactly symmetric.  The parameter $\tauasym$ quantifies
the degree of this asymmetry.

\subsubsection{Price's retrocausality and Wharton's all-at-once picture}

Huw Price has argued on philosophical grounds that the temporal
asymmetry of causation is not fundamental but arises from a
de facto boundary condition---the low-entropy initial state of the
universe~\cite{Price1996}.  In the absence of this boundary
condition, the laws of physics are time-symmetric, and there is no
principled reason to exclude retrocausal influences.
Ken Wharton and collaborators~\cite{Wharton2020} have developed
this into a Lagrangian ``all-at-once'' picture, in which quantum
events are determined by boundary conditions at both ends of a
time interval, rather than by forward-propagating causes.
Leifer and Pusey~\cite{Leifer2017} have shown that, under
reasonable assumptions, time-symmetric ontological models must be
retrocausal.

The $\tauasym$ framework gives these philosophical positions a
quantitative backbone.  Price's argument is, in essence, that
$\tauasym$ \emph{should be} zero at the fundamental level, and
that the observed $\tauasym > 0$ is a consequence of cosmological
boundary conditions (low initial entropy).  The $\tauasym$ parameter
quantifies ``how retrocausal'' a process is: $\tauasym = 0$ is the
perfectly retrocausal (or rather, perfectly a-causal) regime where
Price's symmetry argument holds exactly; $\tauasym = 1$ is the
maximally irreversible regime where retrodiction fails completely.
Wharton's all-at-once picture corresponds to processes for which
the Lagrangian boundary value problem is well-posed---which, in
the Petz language, is precisely the condition $\tauasym = 0$ where
forward and backward descriptions are interchangeable.

The bound~(\ref{eq:tau_bound}) adds a new element:
$\tauasym > 0$ is not arbitrary but bounded by entropy production.
Price's philosophical argument that temporal asymmetry comes from
thermodynamics receives a precise mathematical form:
$\tauasym \leq 1 - e^{-\Sigma/2}$.

\subsubsection{Zurek's decoherence and einselection}

The decoherence program~\cite{Zurek2003,Zeh1970} explains the
emergence of classicality through environment-induced
superselection: the environment monitors the system, selecting
a preferred pointer basis and destroying coherence between
pointer states.  In the modern understanding, decoherence does
not ``collapse'' the wavefunction; it redistributes coherence
from the system to system-environment correlations, making the
coherence locally inaccessible.

In the $\tauasym$ framework, decoherence is precisely the
mechanism that drives $\tauasym$ from zero to a positive value.
Before the system interacts with the environment, the system's
evolution is unitary ($\tauasym = 0$).  After interaction, tracing
out the environment produces entropy ($\Sigma > 0$) and hence
$\tauasym > 0$.  Zurek's einselection is the dynamical process
that creates temporal asymmetry.

This restatement is not merely notational.  It adds three things.
First, $\tauasym$ provides a continuous, quantitative measure of
``how much'' decoherence has occurred, going beyond the binary
distinction between coherent and decohered.
Second, the bound~(\ref{eq:tau_bound}) connects decoherence
directly to thermodynamic entropy production, placing it within
the framework of stochastic thermodynamics.
Third, the equivalence chain~(\ref{eq:equiv_chain}) shows that
decoherence, entropy production, environmental information
leakage, and non-Markovianity are all aspects of the same
underlying fact: $\tauasym > 0$.

\subsubsection{Connes-Rovelli thermal time}

Connes and Rovelli~\cite{Connes1994} proposed that in a generally
covariant theory without a preferred time variable, the flow of
time can be identified with the modular automorphism group of a
thermal state.  The modular flow generated by a state $\omega$
on a von Neumann algebra $\mathcal{M}$ provides a canonical
one-parameter group of automorphisms $\sigma_t^\omega$, and
Connes and Rovelli suggest identifying $t$ with physical time.
In this picture, time is not a fundamental parameter but
emerges from the state.

The connection to $\tauasym$ is suggestive, though we do not
claim it is tight.  In the Connes-Rovelli framework, a system
in a KMS (thermal equilibrium) state has a well-defined modular
flow, and time is the parameter along this flow.  But a KMS
state is a fixed point of its own modular evolution---the system
is in equilibrium, and nothing ``happens.''  The arrow of time
(the distinction between past and future) requires departure
from equilibrium, which in the $\tauasym$ language means
$\Sigma > 0$ and hence $\tauasym > 0$.

The Connes-Rovelli picture and the $\tauasym$ framework agree on
a central point: \emph{when there is no entropy production, there
is no time}.  For a closed system in equilibrium ($\Sigma = 0$),
the modular flow is trivial (no net evolution) and $\tauasym = 0$
(no temporal asymmetry).  Both frameworks locate the origin of
time in irreversibility.  The $\tauasym$ framework adds the
quantitative bound linking the strength of the time arrow to the
amount of entropy produced.

\subsubsection{Page-Wootters mechanism}

Page and Wootters~\cite{Page1983} showed that a static
(zero-energy eigenstate) universe can give rise to the appearance
of time evolution for subsystems.  In their construction, the
total state of a universe composed of a ``clock'' $C$ and a
``rest'' $R$ satisfies $H_{\text{total}}\ket{\Psi} = 0$, so
the total state is time-independent.  But the conditional state
of $R$, given a reading $t$ of the clock, is
$\bra{t}_C\ket{\Psi}_{CR} \propto e^{-iH_R t}\ket{\psi_0}_R$
---it evolves with $t$ as if $t$ were time.

In the $\tauasym$ framework, the Page-Wootters construction is
a special case of $\tauasym = 0$ for the total system.  The
universe is closed; its total evolution is unitary (indeed,
trivially so, since $H\ket{\Psi} = 0$); the Petz recovery is
exact.  There is no temporal asymmetry for the universe as a
whole.  Time---the appearance of directed change---arises only
when we condition on subsystem degrees of freedom (the clock),
which amounts to tracing out part of the system and hence
introducing $\tauasym > 0$ from the subsystem's perspective.

The Page-Wootters mechanism and the $\tauasym$ framework converge
on the same conclusion: \emph{time is the experience of being a
subsystem}.  For the whole, $\tauasym = 0$ and there is no time.
For the part, $\tauasym > 0$ and time flows.

\subsubsection{$\tauasym = 0$ as P-CTC equivalence}

Jean, Silva, and Vilasini~\cite{Jean2025} have recently established
a striking mathematical equivalence: every time-symmetric quantum
process (described by a multi-time state) can be exactly replicated
by a quantum circuit assisted by post-selected closed timelike
curves (P-CTCs, in the sense of Lloyd \emph{et al.}~\cite{Lloyd2011}),
and conversely.  The mapping is constructive and preserves all
measurement statistics.

In the $\tauasym$ language, this yields a sharpened characterization:
\begin{equation}
\tauasym = 0 \;\Longleftrightarrow\; \text{time-symmetric}
\;\Longleftrightarrow\; \text{P-CTC-equivalent}.
\label{eq:pctc}
\end{equation}
A closed system with $\tauasym = 0$ is not merely ``reversible''
in the colloquial sense---it is \emph{operationally equivalent} to
a process with bidirectional information flow, where outputs can
inform inputs through post-selected teleportation.  This is
\emph{not} the claim that closed timelike curves physically exist.
It is the statement that the operational predictions of a
time-symmetric quantum process are indistinguishable from those of
a P-CTC-assisted circuit.

This equivalence illuminates the quantum eraser.  The signal-idler
pair ($\tauasym = 0$) has P-CTC structure: the coincidence counting
that reveals interference is precisely the post-selection step in the
P-CTC protocol.  The ``retrocausal'' appearance of the eraser is not
mysterious---it is the expected behavior of any time-symmetric process.

Crucially, Jean \emph{et al.}\ prove a binary equivalence:
a process either has full P-CTC structure or it does not.
What $\tauasym$ adds is the \emph{continuous interpolation}:
the bound $F \geq \exp(-\Sigma/2)$ quantifies how the P-CTC
structure degrades as the system opens ($\tauasym > 0$).
The transition from $\tauasym = 0$ to $\tauasym > 0$ is the
loss of time-symmetric, P-CTC-equivalent structure---the emergence
of a preferred temporal direction.

\subsubsection{Proper versus improper mixtures: the retrodictive
distinction}

Liu, Bai, and Scarani~\cite{Liu2026} have demonstrated that
the density matrix alone is not sufficient for retrodiction.
Two states with \emph{identical} density matrices---one a
\emph{proper} mixture (representing classical ignorance about which
pure state was prepared) and the other an \emph{improper} mixture
(obtained by tracing out part of an entangled system)---give
\emph{completely different} retrodictions via the Petz map.

Their key example is instructive.  Consider a qubit with
$\beta_S = I/2$ in both cases.  If the prior belief is a proper
mixture (the qubit is $\ket{0}$ or $\ket{1}$ with equal
probability), then observing $\ket{0}$ leads to the retrodiction
$\ket{0}$ with certainty---the observation updates the belief.
If the prior belief is an improper mixture (the qubit is half of a
Bell state), then observing $\ket{0}$ leads to the retrodiction
$I/2$---the observation does not update the belief at all.

This has profound implications for $\tauasym$.  The condition
$\tauasym = 0$ has two qualitatively different meanings:
\begin{enumerate}
\item \emph{Closed system (unitary evolution)}: $F = 1$,
$\Sigma = 0$.  Perfect recovery.  Time is genuinely absent.
\item \emph{Improper mixture (entangled subsystem)}: Recovery is
``trivially perfect'' because the belief never changes.
This is not real reversibility but the \emph{absence of a
definite past}.
\end{enumerate}

The distinction maps directly onto the two regimes of quantum
mechanics:
\begin{itemize}
\item \emph{Improper mixture regime} (``AND''): No retrodiction
is possible---there is no definite past to recover.
Superposition persists.  The system has no retrodictable history.
\item \emph{Proper mixture regime} (``OR''): Retrodiction works
---a definite past exists, and it can be inferred from
the present.  Classical behavior emerges.
\end{itemize}

Decoherence is precisely the transition from improper to proper
mixture.  Liu \emph{et al.}\ show that this transition is
\emph{irreversible at the retrodictive level}: converting an
improper mixture to a proper one changes what inferences are
possible about the past.  This is not merely for-all-practical-
purposes irreversible; it is a qualitative change in the structure
of retrodiction itself.

In the $\tauasym$ framework, the improper$\to$proper transition
is the mechanism by which $\tauasym$ acquires its meaning.  Before
decoherence (improper regime), $\tauasym$ is trivially zero because
there is no past to recover.  After decoherence (proper regime),
$\tauasym > 0$ because a definite past has emerged but is only
partially recoverable.  The arrow of time is not merely the growth
of entropy; it is the emergence of a retrodictable history.

\subsection{What temporal coexistence is not}

Having stated what temporal coexistence means, we must be equally
clear about what it does not mean.  Several potential
misunderstandings must be forestalled.

\emph{It does not mean the future is ``already there.''}
The operational definition~(Definition~\ref{def:coexistence})
says that forward and backward inference are equivalent, not
that the future exists as a spatiotemporal object.  This is a
statement about information, not about ontology.  It is
analogous to saying that a ciphertext and its plaintext contain
the same information (given the key)---this does not mean the
plaintext ``already exists'' somewhere; it means the
transformation between them is lossless.

\emph{It does not violate the second law of thermodynamics.}
The equivalence chain~(\ref{eq:equiv_chain}) shows that
$\tauasym = 0$ if and only if $\Sigma = 0$.  The second law
(for closed systems, $\Sigma = 0$; for open systems,
$\langle\Sigma\rangle \geq 0$ on average) is perfectly
consistent with $\tauasym = 0$ for closed systems.
The point is not that the second law is violated, but that the
second law \emph{is} the arrow of time: $\Sigma > 0$ and
$\tauasym > 0$ are equivalent statements.

\emph{It does not provide a mechanism for backward-in-time
signaling.}
No-signaling is a kinematic constraint that holds independently
of $\tauasym$.  Even when $\tauasym = 0$, the reduced state of
a subsystem is independent of operations performed on a
spacelike-separated partner.  Temporal coexistence means that
the joint state's information content is time-independent, not
that any party can exploit this for communication.

\emph{It does not mean that time is an illusion.}
For any subsystem in an open environment, $\tauasym > 0$ is a
hard physical fact: the past is genuinely not recoverable from
the present, and the arrow of time has measurable consequences.
What the framework suggests is not that time is illusory, but
that it is \emph{emergent}---arising from the subsystem's
relation to its environment rather than from the fundamental
dynamics.

%=======================================================================
\section{The Observer and the Arrow}
\label{sec:observer}
%=======================================================================

\subsection{The observation paradox}

There is a tension in the $\tauasym$ framework that we wish to
state plainly, because it is both a feature and a limitation.

Consider a closed system evolving unitarily.  By
Corollary~(\ref{eq:equiv_chain}), $\tauasym = 0$: the
system has no temporal asymmetry.  Now suppose we wish to
\emph{verify} this.  To do so, we must make a measurement---but
measurement couples the system to an apparatus, breaking the
closure.  The act of observation opens the system, introducing
$\Sigma > 0$ and hence $\tauasym > 0$.

This is not a new observation; it is the quantum measurement
problem in a different guise.  But the $\tauasym$ framework
gives it a precise quantitative form.  Before observation, the
system has $\tauasym = 0$.  After observation, the
system-apparatus composite has $\tauasym > 0$ (from the
system's perspective, since the apparatus has been traced out).
The arrow of time is created by the act of looking.

We do not claim to resolve this paradox.  Rather, we note that
$\tauasym$ provides a continuous parameterization of it.
The transition from $\tauasym = 0$ to $\tauasym > 0$ is not
instantaneous or discontinuous; it occurs gradually as the
system couples to the environment.  The decoherence
timescale~\cite{Zurek2003} is the timescale on which $\tauasym$
grows from zero.  For macroscopic systems, this timescale is
absurdly short (of order $10^{-40}$ seconds for a dust grain in
air), which is why the arrow of time appears so robust.

The Liu-Bai-Scarani result~\cite{Liu2026} sharpens this paradox.
Before observation, the system's mixture is \emph{improper}
(entangled with degrees of freedom the observer has not accessed),
and retrodiction is frozen---the Petz map returns the prior
unchanged.  The act of observation converts the mixture from
improper to proper: the observer gains a definite belief about the
system's state, and retrodiction becomes possible.
But this conversion is itself irreversible: the retrodictive
structure of an improper mixture cannot be recovered from its
proper counterpart.  The observer, by looking, simultaneously
\emph{enables} retrodiction (creating a definite past) and
\emph{destroys} the time-symmetric P-CTC structure
(introducing $\tauasym > 0$).  The price of having a history
is losing the ability to undo it.

\subsection{Subsystem perspective and the emergence of time}

The Page-Wootters mechanism (Section~\ref{sec:coexistence})
suggests a radical conclusion: if the universe as a whole is a
closed system, then $\tauasym_{\text{universe}} = 0$ and there
is no arrow of time for the universe.  Time arises only from the
perspective of subsystems.

Let us make this more precise.  Consider a bipartite system
$AB$ in a pure state $\ket{\Psi}_{AB}$, evolving under a total
Hamiltonian $H_{AB}$.  The total evolution is unitary:
$\ket{\Psi(t)}_{AB} = e^{-iH_{AB}t}\ket{\Psi(0)}_{AB}$.
For the total system, $\tauasym_{\text{total}} = 0$ at all times.

Now consider subsystem $A$ alone.  Its state is
$\rho_A(t) = \Tr_B\big(\op{\Psi(t)}{\Psi(t)}_{AB}\big)$,
and its effective dynamics is given by a time-dependent
completely positive map.  In general, this map is not unitary:
tracing over $B$ produces decoherence and dissipation.
The entropy production of this effective channel satisfies
$\Sigma_A \geq 0$ (on average), and the temporal asymmetry
$\tauasym_A > 0$ whenever $A$ is entangled with $B$.

The arrow of time for subsystem $A$ is therefore a consequence
of its entanglement with $B$.  The more entangled $A$ is with
its environment, the larger $\Sigma_A$, and the stronger the
arrow of time.  This is Zurek's insight~\cite{Zurek2003},
restated in the Petz language.

But the $\tauasym$ framework adds a quantitative dimension.
The bound~(\ref{eq:tau_bound}) says that $\tauasym_A$ is not
merely positive but is bounded by the entropy produced in the
$A$-$B$ interaction.  A strongly coupled subsystem has a strong
arrow of time; a weakly coupled subsystem has a weak one.
An isolated subsystem (if it could exist) would have $\tauasym = 0$
and no arrow of time at all.

This leads to a slogan that captures the framework's content:
\begin{quote}
\emph{Time is the experience of being a subsystem.  The arrow of
time is the cost of tracing out the environment.  The strength of
the arrow is bounded by the entropy produced in the tracing.}
\end{quote}

\subsection{Quantum error correction as artificial closure}

Quantum error correction (QEC)~\cite{KnillLaflamme1997} fits
naturally into this picture.  A quantum error-correcting code
protects a logical subspace from environmental noise by
encoding information redundantly.  The decoder---which, as
established in Huang~\cite{Huang2026}, is approximating the
Petz recovery map---reverses the effects of the noise channel
on the encoded information.

In the $\tauasym$ language, QEC is the technology of
maintaining $\tauasym \approx 0$ against the environment's
constant drive toward $\tauasym > 0$.  The noise channel
$\chan$ pushes $\tauasym$ upward; the decoder $\Petz$ pushes
it back toward zero.  Successful error correction means that
the logical information experiences no temporal asymmetry: the
logical past is perfectly recoverable from the logical present.

This perspective illuminates why QEC is so difficult and so
resource-intensive.  The environment is constantly producing
entropy, driving $\tauasym$ upward.  The decoder must
counteract this entropy production.  The threshold
theorem~\cite{Acharya2025} says that this is possible as long
as the noise rate is below a critical value---but the resources
required scale with the entropy production rate.  In the
$\tauasym$ language: you can maintain $\tauasym_{\text{logical}} = 0$
for the encoded information, but only by accepting a large
$\tauasym_{\text{physical}} > 0$ for the physical qubits and the
classical error correction machinery.

The total entropy production is never negative.  QEC does not
reverse the arrow of time; it redirects it, shielding the logical
information at the expense of increasing entropy elsewhere.  The
analogy to a refrigerator is apt: a refrigerator maintains low
entropy inside by pumping entropy to the environment.  QEC
maintains $\tauasym_{\text{logical}} = 0$ by pumping temporal
asymmetry to the physical layer.

\subsection{The experimenter revisited}

We can now give a complete answer to the question of
Section~\ref{sec:eraser}.  Why must the experimenter wait for
coincidence counting in the quantum eraser?

\begin{enumerate}
\item The entangled signal-idler pair is (approximately) a closed
bipartite system with $\tauasym_{\text{pair}} = 0$.  The
coherence information is present in the joint state and is
time-independent.

\item The experimenter is external to this closed system.  The
experimenter's access to the idler's measurement outcome is
mediated by a chain of irreversible processes: photon detection,
electronic amplification, classical communication, data storage.
Each process produces entropy $\Sigma_i > 0$.

\item By the composition theorem [Huang~\cite{Huang2026},
Theorem~2], the temporal asymmetries compose sub-additively:
$\sqrt{\tauasym_{\text{total}}} \leq \sum_i
\sqrt{\tauasym_i}$.  The experimenter's total temporal
asymmetry is bounded below by the sum of the asymmetries
of the intervening channels.

\item No-signaling guarantees that the signal's marginal
statistics are independent of the idler measurement.  This is
a kinematic constraint that holds regardless of $\tauasym$.

\item The combination of $\tauasym_{\text{exp}} > 0$ and
no-signaling means that the experimenter can access the
correlations only through coincidence counting, which requires
receiving the idler outcome through a time-consuming classical
channel.
\end{enumerate}

The ``waiting'' is therefore the joint consequence of two
independent facts: no-signaling (which prevents direct access
to the correlations through the signal alone) and
$\tauasym > 0$ (which prevents perfect retrodiction of the
joint state from the signal alone).  Neither alone is sufficient.
No-signaling without $\tauasym > 0$ would prevent superluminal
signaling but would allow perfect retrodiction from local
measurements.  $\tauasym > 0$ without no-signaling would
introduce imperfect retrodiction but might allow partial access
to the correlations through the signal's marginal.  Together,
they create the familiar experience of needing to wait for the
classical data from the idler to arrive.

%=======================================================================
\section{Consequences and Predictions}
\label{sec:predictions}
%=======================================================================

\subsection{Testable bounds on recovery fidelity}

The most directly testable prediction of the $\tauasym$ framework
is the recovery bound~(\ref{eq:jrsww}) and its thermodynamic
consequence~(\ref{eq:tau_bound}).  For a quantum channel with
known entropy production $\Sigma$, the bound predicts a
minimum recovery fidelity:
\begin{equation}
F \geq \exp(-\Sigma/2).
\label{eq:testable}
\end{equation}
This can be tested in any platform where both the entropy
production and the Petz recovery fidelity can be measured.

Recent experimental implementations of the Petz recovery map
make this immediately feasible.
Png and Scarani~\cite{Png2025} have implemented single-qubit
Petz recovery on an ion trap quantum processor, achieving
fidelities consistent with theoretical predictions for amplitude
damping, dephasing, and depolarizing channels.
Singh \emph{et al.}~\cite{Singh2025} have implemented the Petz
map on an NMR quantum processor, with similar results.
In both cases, the measured recovery fidelities can be compared
with the entropy production of the noise channel to test the
bound~(\ref{eq:testable}).

A particularly clean test would use amplitude damping
(spontaneous emission, or $T_1$ decay), for which the Petz
recovery achieves exact saturation of the Junge \emph{et al.}
bound for the ground-state input [Huang~\cite{Huang2026},
Theorem~1]: $F^2 = 1/(1+\gamma)$ where $\gamma$ is the damping
parameter.  This is a falsifiable prediction: if the measured
fidelity exceeds $1/(1+\gamma)$, the framework is wrong.

\subsection{Decoder hierarchy and retrodiction fidelity}

Since the Petz map is the unique Bayesian-consistent retrodiction
functor~\cite{Parzygnat2023}, any good quantum error correction
decoder must be approximating it.  The known hierarchy of QEC
decoders---maximum likelihood $>$ neural network
(AlphaQubit~\cite{Bausch2024}) $>$ minimum-weight perfect
matching $>$ Union-Find~\cite{deMartiOlius2024}---should
therefore be ordered by retrodiction fidelity: how well each
decoder approximates the Petz map.

This is a testable prediction.  For each decoder $\mathcal{D}$,
one can compute the channel divergence
$D(\tilde{\Petz}_{\sigma,\chan}\|\mathcal{D})$ and verify that
it correlates with the decoder's logical error rate.  We are not
aware of this comparison having been performed, and we propose it
as a diagnostic tool for decoder design: a decoder's performance
can be benchmarked by its closeness to Petz recovery.

\subsection{Non-Markovian dynamics and transient reversal}

In non-Markovian open quantum systems~\cite{Breuer2016}, the
entropy production $\Sigma$ can be temporarily negative: information
flows back from the environment to the system, transiently
reversing the arrow of time.  In the $\tauasym$ framework, this
corresponds to a transient \emph{decrease} in $\tauasym$---the
system becomes temporarily ``more reversible.''

Formally, for a non-Markovian process described by a time-dependent
map $\chan_t$, the instantaneous entropy production rate
$\dot{\Sigma}(t)$ can be negative during information
backflow~\cite{Breuer2016}.  During these intervals,
the bound~(\ref{eq:tau_bound}) tightens: $\tauasym$ can decrease.
The system is, in a precise operational sense, experiencing a
transient reversal of its arrow of time.

This is not time travel.  The total entropy production
$\Sigma_{\text{total}}$ remains non-negative on average (the
second law holds for the full system-environment composite).
But the subsystem can experience windows of decreasing temporal
asymmetry, during which retrodiction improves.  This is the
$\tauasym$ interpretation of the Breuer-Laine-Piilo measure of
non-Markovianity~\cite{Breuer2016}: non-Markovianity is
precisely the phenomenon of transient $\tauasym$ decrease.

\subsection{The Retrodiction Landauer Principle}

The classical Landauer principle~\cite{Landauer1961} states that
erasing one bit of information costs at least $\kB T \ln 2$ of
heat dissipation.  The $\tauasym$ framework suggests a dual:
\emph{retrodicting} one bit of information about the past of an
irreversible process requires at least $\kB T \ln 2$ of entropy
production to have occurred---otherwise the bit would be
perfectly recoverable and there would be nothing to retrodict.

More precisely: if the entropy production is $\Sigma$, the
maximum retrodiction fidelity is bounded by
$F \geq \exp(-\Sigma/2)$, and the maximum number of
irrecoverable bits is bounded by $\Sigma / \ln 2$.
This ``Retrodiction Landauer Principle'' is the time-reverse of
the standard Landauer principle: where Landauer bounds the cost
of \emph{erasing} information about the present, the retrodiction
principle bounds the cost of \emph{recovering} information about
the past.

Both principles are consequences of the same mathematical
inequality~(\ref{eq:jrsww}).  The standard Landauer bound
reads the inequality in the forward direction (erasure costs
at least this much entropy); the retrodiction bound reads
it in the backward direction (this much entropy production
implies at most this much information loss about the past).

\subsection{Divergence-independent retrodiction for measurements}

A natural objection to the $\tauasym$ framework is that $\tauasym$
depends on the choice of the Uhlmann fidelity (equivalently, the
quantum relative entropy via the JRSWW bound).  Different divergence
measures might yield different retrodiction maps and different
irreversibility parameters.

Kuang, Torii, and Buscemi~\cite{Kuang2025} have recently resolved
this objection for the most physically important class of channels:
\emph{measurement channels} (quantum-to-classical maps).  They prove
that for measurement channels, \emph{all} standard quantum divergences
---Umegaki relative entropy, Petz-R\'enyi, sandwiched R\'enyi,
geometric R\'enyi, Belavkin-Staszewski relative entropy, and trace
distance---yield the \emph{same} retrodiction map: the Petz map.
The unique optimizer is
\begin{equation}
\hat{\gamma}_{Q|x} = \frac{\sqrt{\gamma_Q}\, M_x\, \sqrt{\gamma_Q}}
{\Tr\{M_x \gamma_Q\}},
\label{eq:kuang}
\end{equation}
which is precisely the Petz transpose output.

The mathematical reason is that the semiclassical (direct-sum)
structure of measurement channels allows the divergence to
decompose into independent conditional terms, each of which is
minimized by the same state regardless of which divergence is used.
For general quantum-to-quantum channels (where the output retains
coherence), this uniqueness breaks down: different divergences give
different retrodictions.

For the $\tauasym$ framework, this result is significant:
\emph{$\tauasym = 1 - F_{\text{Petz}}$ is not one choice among
many for quantifying measurement irreversibility---it is the unique
one}.  The bound $F \geq \exp(-\Sigma/2)$ is not contingent on an
arbitrary choice of divergence; it is the bound.  This strengthens
every prediction in Section~\ref{sec:predictions} that involves
measurement channels.

\subsection{Collapse as $\tauasym \to 1$: the emergence of history}

The $\tauasym$ framework provides a natural language for discussing
wavefunction collapse.  When $\tauasym = 0$, the system's state
is perfectly recoverable: no information about the pre-measurement
state has been lost.  When $\tauasym \to 1$, recovery fails
completely: the post-measurement state contains no information
about the pre-measurement state.  This is operationally
indistinguishable from ``collapse.''

The Liu-Bai-Scarani result~\cite{Liu2026} enriches this picture.
Before measurement, the system-apparatus state is an improper
mixture: no definite past can be assigned, and retrodiction is
frozen.  After measurement, the mixture becomes proper: a definite
past exists, and the Petz map can retrodict it.  The measurement
transition is not merely an increase in $\tauasym$---it is a
qualitative change in what retrodiction \emph{means}:
\begin{itemize}
\item \emph{Before} (improper, ``AND''): The system is in
state $A$ \emph{and} state $B$ simultaneously.  No retrodiction
is possible.  $\tauasym$ is trivially zero.
\item \emph{After} (proper, ``OR''): The system was in state $A$
\emph{or} state $B$.  Retrodiction is possible but imperfect.
$\tauasym > 0$ reflects the cost of the conversion.
\end{itemize}

The AND$\to$OR transition IS collapse, operationally defined.
The Kuang-Torii-Buscemi result~\cite{Kuang2025} ensures that this
transition is diagnosed identically by all divergence measures:
the emergence of a definite past is not an artifact of our choice
of metric.

This does not solve the measurement problem---it does not explain
why or when the AND$\to$OR transition occurs, or whether it is
objective or perspectival.  But it does provide a quantitative
framework for discussing collapse that is continuous, bounded by
thermodynamics, divergence-independent for measurements, and
agnostic about interpretation.

Galley, Giacomini, and Selby~\cite{Galley2023} have shown that
coupling a quantum system to a classical gravitational field
necessarily introduces irreversibility---the quantum channel
describing the interaction is not unitary.  In the $\tauasym$
framework, this means that classical gravity contributes to
$\tauasym > 0$.  Whether this contribution is sufficient to
explain the observed rate of collapse (as in the Penrose-Di\'osi
proposal) is an open question.
Danielson, Satishchandran, and Wald~\cite{Danielson2022} have
further shown that any system that radiates (e.g., Bremsstrahlung
from a superposed charge) necessarily decoheres, establishing a
fundamental link between radiation, information loss, and
$\tauasym > 0$.

%=======================================================================
\section{Limitations and Open Questions}
\label{sec:limitations}
%=======================================================================

We have argued that the temporal asymmetry parameter $\tauasym$
provides a quantitative explanation for why the experimenter must
wait in the quantum eraser, and that this connects to the arrow
of time, decoherence, quantum error correction, and stochastic
thermodynamics.  We now turn to the limitations of this framework
and the questions it leaves open.

\subsection{The observation paradox, restated precisely}

The most fundamental limitation is the observation paradox stated
in Section~\ref{sec:observer}.  To verify that $\tauasym = 0$
for a closed system, we must measure it, which opens the system
and makes $\tauasym > 0$.  The condition $\tauasym = 0$ is
therefore verifiable only indirectly: we can verify that the Petz
recovery achieves high fidelity (by preparing many copies,
applying the channel, applying the recovery, and performing
state tomography), and we can verify that $\Sigma \approx 0$
(by carefully controlling the environment), but we cannot verify
$\tauasym = 0$ \emph{without} introducing some $\tauasym > 0$
through the measurement process.

This is not unique to the $\tauasym$ framework; it is a general
feature of quantum mechanics.  But it does mean that the
interpretive claim---that $\tauasym = 0$ implies ``temporal
coexistence''---is not directly falsifiable.  What is falsifiable
is the \emph{bound}: the prediction that $F \geq \exp(-\Sigma/2)$
for any process with entropy production $\Sigma$.  If this bound
is violated, the framework is wrong.  If the bound is satisfied,
the framework is consistent, but the interpretation remains a
choice.

\subsection{Falsifiability}

Let us be precise about what is and is not falsifiable.

\emph{Falsifiable:}
\begin{itemize}
\item The recovery bound $F^2 \geq \exp(-\Delta D)$
[Junge \emph{et al.}~\cite{Junge2018}].  This is a mathematical
theorem, but its physical content is testable: given a measured
$\Delta D$ and a measured $F$, the inequality must hold.
\item The saturation condition
[Huang~\cite{Huang2026}, Theorem~1]:
$F^2 = \exp(-\Delta D)$ if and only if the channel outputs
commute and have constant likelihood ratio.  This predicts exact
fidelity values for specific channels and inputs.
\item The composition inequality
[Huang~\cite{Huang2026}, Theorem~2]:
$\sqrt{\tauasym_{12}} \leq \sqrt{\tauasym_1} +
\sqrt{\tauasym_2^{\text{eff}}}$.  This can be tested by
measuring $\tauasym$ for sequential channels.
\item The decoder hierarchy prediction: better decoders should
approximate the Petz map more closely.
\end{itemize}

\emph{Partially falsifiable:}
\begin{itemize}
\item The equivalence chain~(\ref{eq:equiv_chain}).  Each
individual equivalence is a mathematical theorem, but the
physical interpretation (``$\tauasym = 0$ means no arrow of
time'') is a definition, not a prediction.
\item The Retrodiction Landauer Principle.  It is a restatement
of the recovery bound, so it is falsifiable to the same degree;
but the name ``Retrodiction Landauer'' is an interpretive label.
\end{itemize}

\emph{Not falsifiable (interpretive):}
\begin{itemize}
\item The claim that $\tauasym = 0$ implies ``temporal
coexistence.''  This is a definition
(Definition~\ref{def:coexistence}), not a prediction.  It can be
adopted or rejected without changing any experimental prediction.
\item The slogan ``time is the experience of being a subsystem.''
This is a philosophical interpretation of the mathematical
framework, not a testable hypothesis.
\end{itemize}

We regard this hierarchy of falsifiability as a strength, not a
weakness.  The mathematical content is rigorous and testable;
the interpretation is offered as a coherent way of reading the
mathematics, not as a dogma.

\subsection{Free will and agency}

If $\tauasym = 0$ for a closed system implies that past and future
states contain the same information, does this imply determinism?
Does it eliminate free will?

This question requires care.  The condition $\tauasym = 0$ applies
to \emph{closed} systems---systems that do not interact with
anything external.  Any agent capable of making choices must be
an open system: the agent receives information from the environment
(perception), processes it (cognition), and acts on the
environment (action).  Each of these steps involves irreversible
information processing with $\Sigma > 0$.
An agent, almost by definition, has $\tauasym > 0$.

The $\tauasym$ framework is therefore compatible with libertarian
free will (the agent's choices are not determined by the past,
because the agent's $\tauasym > 0$ means the past is not
perfectly recoverable from the present) and with compatibilist
free will (the agent acts on reasons, and the irreversibility of
the agent's information processing is what makes the action
``free'' in the compatibilist sense).  It is incompatible with
the claim that an agent can be a closed system---but this seems
to be an analytic truth about agency, not a limitation of the
framework.

\subsection{What the framework does not do}

We wish to be explicit about several things the $\tauasym$
framework does \emph{not} accomplish.

\emph{It does not solve the measurement problem.}
The framework provides a continuous parameterization of
decoherence ($\tauasym$ increasing from 0 toward 1) but does not
explain why measurements have definite outcomes.  The transition
from $\tauasym \approx 0$ (coherent superposition) to
$\tauasym \approx 1$ (definite outcome) is described but not
explained.  The ``hard'' measurement problem---why there is a
definite outcome at all, rather than a superposition of
observer-states---is untouched.

\emph{It does not provide a theory of quantum gravity.}
The gravitational sector of $\tauasym$---what happens when
$\tauasym$ is applied to systems in curved spacetime---is
unexplored in this paper.  The suggestive connection between
$\tauasym$ and black hole entropy (where information loss is
maximized and $\tauasym \to 1$) deserves investigation but is
beyond our scope.  We note that the Petz map appears in
holographic entanglement wedge reconstruction~\cite{Chen2020},
suggesting that the $\tauasym$ framework may have gravitational
implications, but we do not pursue this here.

\emph{It does not explain the cosmological arrow of time.}
The framework explains why subsystems experience an arrow of time
(because $\Sigma > 0$ for open systems), but it does not explain
why the initial state of the universe had low entropy.  The
cosmological arrow is an input to the framework (it determines
the magnitude of $\Sigma$ and hence of $\tauasym$), not an output.
This is a limitation shared with all approaches that derive the
thermodynamic arrow from boundary conditions, including
Price's~\cite{Price1996}.

\subsection{Open questions}

Several questions arise naturally from this framework and deserve
further investigation.

\emph{Gravitational $\tauasym$:}
Can $\tauasym$ be defined for gravitational degrees of freedom?
The Bekenstein-Hawking entropy suggests $\tauasym \to 1$ at a
black hole horizon (maximal information loss), while flat
spacetime should have $\tauasym = 0$ (no gravitational
irreversibility).  A gravitational $\tauasym$ would need to be
defined in terms of a quantum channel acting on matter crossing
the horizon.

\emph{Many-body scaling:}
How does $\tauasym$ scale with system size?  For $n$ qubits
interacting with an environment, is $\tauasym$ extensive,
intensive, or something else?  The composition
theorem~\cite{Huang2026} gives bounds on sequential channels,
but the scaling for parallel channels (many qubits simultaneously
decohering) is unknown.

\emph{Cosmological arrow:}
The $\tauasym$ framework derives the thermodynamic arrow from
entropy production but does not explain the cosmological arrow.
Can a Page-Wootters-type construction, combined with the
$\tauasym$ framework, provide a self-consistent account of the
cosmological arrow without postulating low initial entropy?

\emph{Collapse dynamics:}
Does the rate of $\tauasym$ increase during measurement
correspond to any known collapse model (GRW, Di\'osi-Penrose,
CSL)?  If so, the $\tauasym$ framework could provide a
thermodynamic perspective on objective collapse theories,
bounding the collapse rate by entropy production.

\emph{Non-Markovian structure:}
The equivalence chain~(\ref{eq:equiv_chain}) is stated for
Markovian channels.  Can it be extended to non-Markovian
processes using the process tensor
framework~\cite{Pollock2018}?  If so, the temporal
coexistence condition $\tauasym = 0$ would need to be
reformulated for multi-time processes.

%=======================================================================
\section{Conclusion}
\label{sec:conclusion}
%=======================================================================

The delayed-choice quantum eraser poses a simple question:
why must the experimenter wait?  The standard answer---no-signaling
prevents superluminal communication---is correct but incomplete.
It identifies the constraint but not the mechanism.

We have argued that the mechanism is the experimenter's nonzero
temporal asymmetry $\tauasym$, arising from the irreversible
classical channels through which the experimenter accesses the
entangled pair.  The entangled pair itself has $\tauasym = 0$:
for it, past and future states contain the same information, and
the Petz recovery map---the unique Bayesian-consistent retrodiction
functor---achieves perfect fidelity.  The experimenter's
$\tauasym > 0$ is bounded by the entropy production in the
measurement and communication chain, via
$\tauasym \leq 1 - \exp(-\Sigma/2)$.

This interpretation connects to a broader picture of the arrow of
time.  Time, in this framework, is not a fundamental feature of
dynamics but the operational cost of being a subsystem in an open
environment.  For the universe as a whole (if it is closed),
$\tauasym = 0$ and there is no arrow of time.  For subsystems,
$\tauasym > 0$ because tracing out the environment produces entropy.
The strength of the arrow is bounded by the amount of entropy
produced.

Recent results sharpen this picture in three ways.
Jean, Silva, and Vilasini~\cite{Jean2025} show that $\tauasym = 0$
is operationally equivalent to having access to post-selected closed
timelike curves---giving ``temporal coexistence'' a precise
information-theoretic meaning.
Liu, Bai, and Scarani~\cite{Liu2026} show that
the transition from $\tauasym = 0$ to $\tauasym > 0$ is the
transition from improper to proper mixture---from ``no definite
past'' to ``definite but partially recoverable past.''
Kuang, Torii, and Buscemi~\cite{Kuang2025} show that for
measurement channels, $\tauasym$ is the unique irreversibility
measure: all divergences yield the same retrodiction.

Together with the existing connections to Zurek's decoherence program,
Aharonov's two-state vector formalism, Price's retrocausality,
Connes-Rovelli thermal time, and the Page-Wootters mechanism, these
results support a unified picture: the arrow of time is the emergence
of a retrodictable history from a time-symmetric substrate.

The framework makes testable predictions: recovery fidelity bounds,
saturation conditions for specific channels, composition
inequalities, and decoder hierarchy orderings.  It also makes
interpretive claims that are not directly falsifiable but are
offered as a coherent reading of the mathematics.

We end with a reflection.  The quantum eraser is often presented
as a paradox about the nature of time---as evidence that the
future can influence the past, or that quantum mechanics is
incompatible with our everyday notion of temporal order.  The
$\tauasym$ framework dissolves this paradox by locating the source
of temporal order not in the fundamental dynamics (which are
time-symmetric) but in the observer's relation to the observed
system.  The entangled pair has no arrow of time.  The experimenter
does.  The gap between them is what we call ``waiting.''

\subsection*{Acknowledgments}
The author thanks M.~M.~Wilde for valuable feedback on earlier work,
and A.~J.~Parzygnat and F.~Buscemi for foundational contributions to
the theory of quantum retrodiction.  This work builds on the framework
established in Ref.~\cite{Huang2026}.

%=======================================================================
% BIBLIOGRAPHY
%=======================================================================

\begin{thebibliography}{50}

\bibitem{Scully1982}
M.~O. Scully and K.~Dr\"uhl,
``Quantum eraser: A proposed photon correlation experiment concerning
observation and `delayed choice' in quantum mechanics,''
Phys.\ Rev.\ A \textbf{25}, 2208 (1982).

\bibitem{Kim2000}
Y.-H. Kim, R.~Yu, S.~P. Kulik, Y.~Shih, and M.~O. Scully,
``Delayed `choice' quantum eraser,''
Phys.\ Rev.\ Lett.\ \textbf{84}, 1 (2000).

\bibitem{Petz1986}
D.~Petz,
``Sufficient subalgebras and the relative entropy of states of a
von Neumann algebra,''
Commun.\ Math.\ Phys.\ \textbf{105}, 123 (1986).

\bibitem{Petz1988}
D.~Petz,
``Sufficiency of channels over von Neumann algebras,''
Q.\ J.\ Math.\ \textbf{39}, 97 (1988).

\bibitem{Parzygnat2023}
A.~J. Parzygnat and F.~Buscemi,
``Axioms for retrodiction: Achieving time-reversal symmetry with a
prior,''
Quantum \textbf{7}, 1013 (2023).

\bibitem{Junge2018}
M.~Junge, R.~Renner, D.~Sutter, M.~M. Wilde, and A.~Winter,
``Universal recovery maps and approximate sufficiency of quantum
relative entropy,''
Ann.\ Henri Poincar\'e \textbf{19}, 2955 (2018).

\bibitem{Fawzi2015}
O.~Fawzi and R.~Renner,
``Quantum conditional mutual information and approximate Markov
chains,''
Commun.\ Math.\ Phys.\ \textbf{340}, 575 (2015).

\bibitem{Kwon2019}
H.~Kwon and M.~S. Kim,
``Fluctuation theorems for a quantum channel,''
Phys.\ Rev.\ X \textbf{9}, 031029 (2019).

\bibitem{Huang2026}
S.-K. Huang,
``The arrow of time from Petz recovery,''
arXiv:2603.XXXXX (2026).

\bibitem{Barnum2002}
H.~Barnum and E.~Knill,
``Reversing quantum dynamics with near-optimal quantum and classical
fidelity,''
J.\ Math.\ Phys.\ \textbf{43}, 2097 (2002).

\bibitem{Crooks1999}
G.~E. Crooks,
``Entropy production fluctuation theorem and the nonequilibrium
work relation for free energy differences,''
Phys.\ Rev.\ E \textbf{60}, 2721 (1999).

\bibitem{Kolchinsky2017}
A.~Kolchinsky and D.~H. Wolpert,
``Dependence of dissipation on the initial distribution over states,''
J.\ Stat.\ Mech.\ \textbf{2017}, 083202 (2017).

\bibitem{ReebWolf2014}
D.~Reeb and M.~M. Wolf,
``An improved Landauer principle with finite-size corrections,''
New J.\ Phys.\ \textbf{16}, 103011 (2014).

\bibitem{ABL1964}
Y.~Aharonov, P.~G. Bergmann, and J.~L. Lebowitz,
``Time symmetry in the quantum process of measurement,''
Phys.\ Rev.\ \textbf{134}, B1410 (1964).

\bibitem{Aharonov1990}
Y.~Aharonov and L.~Vaidman,
``Properties of a quantum system during the time interval between
two measurements,''
Phys.\ Rev.\ A \textbf{41}, 11 (1990).

\bibitem{Price1996}
H.~Price,
\emph{Time's Arrow and Archimedes' Point: New Directions for the
Physics of Time}
(Oxford University Press, Oxford, 1996).

\bibitem{Wharton2020}
K.~B. Wharton and N.~Argaman,
``Colloquium: Bell's theorem and locally mediated reformulations
of quantum mechanics,''
Rev.\ Mod.\ Phys.\ \textbf{92}, 021002 (2020).

\bibitem{Leifer2017}
M.~S. Leifer and M.~F. Pusey,
``Is a time symmetric interpretation of quantum theory possible
without retrocausality?''
Proc.\ R.\ Soc.\ A \textbf{473}, 20160607 (2017).

\bibitem{Zurek2003}
W.~H. Zurek,
``Decoherence, einselection, and the quantum origins of the
classical,''
Rev.\ Mod.\ Phys.\ \textbf{75}, 715 (2003).

\bibitem{Zeh1970}
H.~D. Zeh,
``On the interpretation of measurement in quantum theory,''
Found.\ Phys.\ \textbf{1}, 69 (1970).

\bibitem{Connes1994}
A.~Connes and C.~Rovelli,
``Von Neumann algebra automorphisms and time-thermodynamics
relation in generally covariant quantum theories,''
Class.\ Quantum Grav.\ \textbf{11}, 2899 (1994).

\bibitem{Page1983}
D.~N. Page and W.~K. Wootters,
``Evolution without evolution: Dynamics described by stationary
observables,''
Phys.\ Rev.\ D \textbf{27}, 2885 (1983).

\bibitem{Maccone2009}
L.~Maccone,
``Quantum solution to the arrow-of-time dilemma,''
Phys.\ Rev.\ Lett.\ \textbf{103}, 080401 (2009).

\bibitem{Watanabe1955}
S.~Watanabe,
``Symmetry of physical laws. Part III. Prediction and retrodiction,''
Rev.\ Mod.\ Phys.\ \textbf{27}, 179 (1955).

\bibitem{Ghirardi1980}
G.~C. Ghirardi, A.~Rimini, and T.~Weber,
``A general argument against superluminal transmission through the
quantum mechanical measurement process,''
Lett.\ Nuovo Cim.\ \textbf{27}, 293 (1980).

\bibitem{KnillLaflamme1997}
E.~Knill and R.~Laflamme,
``Theory of quantum error-correcting codes,''
Phys.\ Rev.\ A \textbf{55}, 900 (1997).

\bibitem{Landauer1961}
R.~Landauer,
``Irreversibility and heat generation in the computing process,''
IBM J.\ Res.\ Dev.\ \textbf{5}, 183 (1961).

\bibitem{Landi2021}
G.~T. Landi and M.~Paternostro,
``Irreversible entropy production: From classical to quantum,''
Rev.\ Mod.\ Phys.\ \textbf{93}, 035008 (2021).

\bibitem{Breuer2016}
H.-P. Breuer, E.~M. Laine, J.~Piilo, and B.~Vacchini,
``Colloquium: Non-Markovian dynamics in open quantum systems,''
Rev.\ Mod.\ Phys.\ \textbf{88}, 021002 (2016).

\bibitem{Buscemi2024}
G.~Bai, F.~Buscemi, and V.~Scarani,
``Fully quantum stochastic entropy production,''
arXiv:2412.12489 (2024).

\bibitem{Png2025}
W.-H. Png and V.~Scarani,
``Petz recovery maps of single-qubit decoherence channels in an
ion trap quantum processor,''
Phys.\ Rev.\ A \textbf{112}, 022613 (2025).

\bibitem{Singh2025}
G.~Singh, R.~S. Sahani, V.~Jagadish, L.~Lautenbacher,
N.~K. Bernardes, and K.~Dorai,
``Realizing the Petz recovery map on an NMR quantum processor,''
arXiv:2508.08998 (2025).

\bibitem{Chen2020}
C.-F. Chen, G.~Penington, and G.~Salton,
``Entanglement wedge reconstruction using the Petz map,''
J.\ High Energy Phys.\ \textbf{01}, 168 (2020).

\bibitem{Pollock2018}
F.~A. Pollock, C.~Rodr\'iguez-Rosario, T.~Frauenheim,
M.~Paternostro, and K.~Modi,
``Non-Markovian quantum processes: Complete framework and efficient
characterization,''
Phys.\ Rev.\ A \textbf{97}, 012127 (2018).

\bibitem{Galley2023}
T.~D. Galley, F.~Giacomini, and J.~H. Selby,
``Any consistent coupling between classical gravity and quantum
matter is fundamentally irreversible,''
Quantum \textbf{7}, 1142 (2023).

\bibitem{Danielson2022}
D.~L. Danielson, G.~Satishchandran, and R.~M. Wald,
``Gravitationally mediated entanglement: Newtonian field versus
gravitons,''
Phys.\ Rev.\ D \textbf{105}, 086001 (2022).

\bibitem{Acharya2025}
R.~Acharya \emph{et al.} (Google Quantum AI),
``Quantum error correction below the surface code threshold,''
Nature \textbf{638}, 920 (2025).

\bibitem{Bausch2024}
J.~Bausch, A.~W. Senior \emph{et al.},
``Learning high-accuracy error decoding for quantum processors,''
Nature \textbf{635}, 834 (2024).

\bibitem{deMartiOlius2024}
A.~deMarti~iOlius \emph{et~al.},
``Decoding algorithms for surface codes,''
Quantum \textbf{8}, 1498 (2024).

\bibitem{Fullwood2023}
A.~J. Parzygnat and J.~Fullwood,
``From time-reversal symmetry to quantum Bayes' rules,''
PRX Quantum \textbf{4}, 020334 (2023).

\bibitem{Frank2025}
R.~L. Frank,
``Quantum thermodynamic semidefinite programs,''
arXiv:2508.19301 (2025).

\bibitem{Wheeler1978}
J.~A. Wheeler,
``The `past' and the `delayed-choice' double-slit experiment,''
in \emph{Mathematical Foundations of Quantum Theory},
edited by A.~R. Marlow (Academic Press, 1978), pp.~9--48.

\bibitem{Aw2024}
C.~C. Aw, L.~H. Zaw, M.~Balanzo-Juando, and V.~Scarani,
``Role of dilations in reversing physical processes: Tabletop
reversibility and generalized thermal operations,''
PRX Quantum \textbf{5}, 010332 (2024).

\bibitem{Jarzynski1997}
C.~Jarzynski,
``Nonequilibrium equality for free energy differences,''
Phys.\ Rev.\ Lett.\ \textbf{78}, 2690 (1997).

\bibitem{Sagawa2012}
T.~Sagawa and M.~Ueda,
``Fluctuation theorem with information exchange,''
Phys.\ Rev.\ Lett.\ \textbf{109}, 180602 (2012).

\bibitem{Jean2025}
E.~Jean, R.~Silva, and V.~Vilasini,
``An equivalence between time-symmetry and cyclic causality in
quantum theory,''
arXiv:2508.02463 (2025).

\bibitem{Lloyd2011}
S.~Lloyd, L.~Maccone, R.~Garcia-Patron, V.~Giovannetti, and
Y.~Shikano,
``Quantum mechanics of time travel through post-selected
teleportation,''
Phys.\ Rev.\ D \textbf{84}, 025007 (2011).

\bibitem{Liu2026}
M.~Liu, G.~Bai, and V.~Scarani,
``Proper and improper mixed states serve as different prior beliefs
for quantum state retrodiction,''
Phys.\ Rev.\ Lett.\ \textbf{136}, 060203 (2026).

\bibitem{Kuang2025}
J.~Kuang, K.~Torii, and F.~Buscemi,
``Quantum measurement retrodiction and entropic uncertainty
relations,''
arXiv:2511.20281 (2025).

\bibitem{BaiBuscemiScarani2025}
G.~Bai, F.~Buscemi, and V.~Scarani,
``The Petz recovery map as the optimal retrodiction map,''
Phys.\ Rev.\ Lett.\ \textbf{135}, 200201 (2025).

\end{thebibliography}

\end{document}
